
\subsubsection{Supervised-Unsupervised Performance Gap}

The central finding is a 77-percentage-point gap between supervised probe accuracy (79.6\%) and unsupervised clustering performance (2.3\% NMI normalized to percentage scale). This gap indicates that lifecycle categories are encoded in semantic embeddings as signals accessible to supervised methods but do not manifest as natural clusters detectable by standard unsupervised algorithms.

Multiple lines of evidence support signal existence: inter-annotator agreement ($\kappa$=0.645 in adjusted methodology), linear probe accuracy (79.6\% overall, 97-100\% in repository-specific best cases), and consistent cross-repository supervised generalization (2.9pp standard deviation). The signal exists but requires task-specific amplification to access.

\subsubsection{Attribution to Class Imbalance}

The clustering failure appears attributable to severe class imbalance (8.3\% Responsible AI, 11:1 ratio). K-means assumes balanced, well-separated clusters; extreme imbalance causes the minority class to appear as noise. Evidence:

\begin{enumerate}
\item Cluster sizes (287:13) nearly match the ground truth distribution (275:25), indicating trivial majority-class assignment rather than structure discovery.
\item Repository variance correlates with class balance: UCI (14\% RAI, NMI 0.394) versus HuggingFace (5.3\% RAI, NMI 0.013).
\item Supervised probes overcome imbalance through learned class weights, achieving 77.8\% RAI recall; K-means achieves approximately 52\% (13/25), barely above chance.
\end{enumerate}

This evidence is correlational; controlled rebalancing experiments were not performed. However, the pattern is consistent across multiple analyses.

\subsubsection{Alternative Explanations}

\textbf{Repository heterogeneity}: Differences in schema structure (UCI plain text versus HuggingFace YAML) may contribute to clustering variance beyond class balance. However, supervised probes generalize consistently (91-96\% accuracy), suggesting heterogeneity affects unsupervised methods specifically.

\textbf{Vocabulary variation}: The lexical baseline achieved 0\% RAI recall, demonstrating extreme terminology diversity. However, semantic embeddings successfully encode this variation for supervised detection (77.8\% recall), indicating vocabulary is not a fundamental barrier when supervision is available.

\textbf{Schema complexity}: Structured repositories may introduce artifacts. However, unscaffolded HuggingFace data achieved 94.7\% probe accuracy, validating signal existence independent of explicit templates.

\subsubsection{Limitations}

\textbf{Scale}: The 300-sample proof-of-concept demonstrates signal existence but does not validate production-scale deployment. Larger repositories (Kaggle, Zenodo) and additional sample sizes are necessary to confirm scalability.

\textbf{Annotation method}: Inter-annotator agreement was simulated via content-based heuristics rather than measured with human expert annotators. The ground truth file reports an adjusted $\kappa$ of 0.645, but this requires independent validation with actual human coders.

\textbf{Causal evidence}: The attribution to class imbalance is supported by correlational evidence but lacks controlled experimental manipulation (e.g., artificial rebalancing).

\textbf{Temporal generalization}: Metadata conventions evolve. The 2026 snapshot may not generalize to future documentation practices.

\textbf{Production accuracy requirements}: The 79.6\% probe accuracy includes a 22\% false negative rate (4 of 18 RAI fields missed). Whether this suffices for production applications depends on use-case requirements.

\subsubsection{Deployment Implications}

The results indicate that unsupervised cross-repository lifecycle detection is not feasible under current metadata distributions. However, supervised and semi-supervised approaches remain viable:

\begin{enumerate}
\item \textbf{Few-shot learning}: 60 training samples achieved 76\% accuracy on a held-out repository. Optimization may reduce annotation requirements to 10-20 samples per repository, though this requires validation.
\end{enumerate}

\begin{enumerate}
\item \textbf{Active learning}: Select maximally informative samples to minimize annotation costs, targeting high-variance fields and underrepresented repositories.
\end{enumerate}

\begin{enumerate}
\item \textbf{Probe-based label propagation}: Train probes on small labeled sets and propagate predictions. Cross-repository consistency (2.9pp standard deviation) suggests probes generalize with minimal adaptation.
\end{enumerate}

\begin{enumerate}
\item \textbf{Human-in-the-loop validation}: For critical applications (e.g., flagging ethical concerns), the 22\% false negative rate may be unacceptable, requiring human review of predictions.
\end{enumerate}

Production deployment requires addressing unresolved questions: Does 79.6\% accuracy suffice for intended applications? How do annotation costs scale to thousands of datasets? Can active learning reduce budgets sufficiently for ecosystem-wide deployment?

